
% Version 2.6: settlement-consistent EV--BESS mathematical formulation.
\documentclass[final,5p,times]{elsarticle}

\usepackage[utf8]{inputenc}
\usepackage[T1]{fontenc}
\usepackage{amsmath,amssymb,amsfonts}
\usepackage{booktabs}
\usepackage{graphicx}
\usepackage{xcolor}
\usepackage{url}
% \usepackage{hyperref} % disabled for cleaner frontmatter/bookmark compilation
\IfFileExists{stfloats.sty}{\usepackage{stfloats}}{}
\usepackage{array}
\usepackage{tabularx}
\usepackage{multirow}
\usepackage{enumitem}
\usepackage{makecell}
\usepackage{algorithm}
\usepackage{algorithmic}

\journal{Sustainable Energy, Grids and Networks}
\graphicspath{{figures/}{./}}
\biboptions{sort&compress}

\newcommand{\DeltaT}{\Delta t}
\newcommand{\EV}{\mathrm{EV}}
\newcommand{\BESS}{\mathrm{BESS}}
\newcommand{\DA}{\mathrm{DA}}
\newcommand{\RT}{\mathrm{RT}}
\newcommand{\GI}{\mathrm{GI}}
\newcommand{\PV}{\mathrm{PV}}
\newcommand{\load}{\mathrm{load}}
\newcommand{\TOU}{\mathrm{TOU}}
\newcommand{\PD}{\mathrm{PD}}
\newcommand{\NCD}{\mathrm{NCD}}
\newcommand{\SOC}{\mathrm{SOC}}
\newcommand{\soc}{\mathrm{soc}}
\newcommand{\RU}{\mathrm{RU}}
\newcommand{\RD}{\mathrm{RD}}
\newcommand{\SP}{\mathrm{SP}}
\newcommand{\NSP}{\mathrm{NSP}}
\newcommand{\WM}{\mathrm{WM}}
\newcommand{\LMP}{\mathrm{LMP}}
\newcommand{\NWM}{\mathrm{NWM}}
\newcommand{\ch}{\mathrm{ch}}
\newcommand{\dch}{\mathrm{dch}}
\newcommand{\up}{\mathrm{up}}
\newcommand{\down}{\mathrm{down}}
\newcommand{\tot}{\mathrm{tot}}
\newcommand{\base}{\mathrm{base}}
\newcommand{\forecast}{\mathrm{forecast}}
\newcommand{\real}{\mathrm{real}}
\newcommand{\executed}{\mathrm{executed}}
\newcommand{\req}{\mathrm{req}}
\newcommand{\net}{\mathrm{net}}
\newcommand{\energy}{\mathrm{energy}}
\newcommand{\capacity}{\mathrm{capacity}}
\newcommand{\avail}{\mathrm{avail}}
\newcommand{\impl}{\mathrm{impl}}
\newcommand{\deploy}{\mathrm{deploy}}
\newcommand{\pen}{\mathrm{pen}}
\newcommand{\calH}{\mathcal{H}}
\newcommand{\calT}{\mathcal{T}}
\newcommand{\calV}{\mathcal{V}}
\newcommand{\calK}{\mathcal{K}}
\newcommand{\calS}{\mathcal{S}}
\newcommand{\calD}{\mathcal{D}}
\newcommand{\calI}{\mathcal{I}}
\newcommand{\pos}[1]{\left[#1\right]_+}
\providecommand{\sep}{; }

% Safe figure command: the manuscript still compiles while figures are being regenerated.
\newcommand{\MaybeFigure}[2]{%
\IfFileExists{#1}{\includegraphics[#2]{#1}}{%
\fbox{\parbox[c][0.19\textheight][c]{0.86\linewidth}{\centering\small Missing figure file.\\ Replace with the paper-quality figure generated from the plotting notebook.}}}}

\begin{document}

\begin{frontmatter}

\title{Execution-Aware Retail--Wholesale Value Stacking for Campus EV Charging and Battery Storage}

\author[ucsd]{Yizhan Gu}
\ead{yizhangu@ucsd.edu}
\author[totalenergies]{Wente Zeng}
\ead{wente.zeng@totalenergies.com}
\author[totalenergies]{Hadi Eshraghi}
\ead{hadi.eshraghi@external.totalenergies.com}
\author[ucsd]{Jan Kleissl\corref{cor1}}
\ead{jkleissl@ucsd.edu}
\cortext[cor1]{Corresponding author.}

\affiliation[ucsd]{organization={Center for Energy Research, University of California San Diego},
            city={La Jolla},
            postcode={92093},
            state={CA},
            country={USA}}

\affiliation[totalenergies]{organization={TotalEnergies},
            city={Houston},
            state={TX},
            country={USA}}

\begin{highlights}
\item Execution-aware retail--wholesale value stacking is evaluated for campus EV charging and BESS.
\item Shrinking and rolling MPC are compared under perfect and persistence EV information.
\item Unidirectional EV flexibility and bidirectional BESS operation share one BTM feasibility envelope.
\item A DA/RT two-settlement representation preserves DA energy commitments during RT execution.
\item Net value is measured from executed dispatch after retail tariff, EV-service, and market interactions.
\end{highlights}

\begin{abstract}
Campus electric-vehicle (EV) charging sites can become substantial behind-the-meter (BTM) distributed-energy-resource portfolios when coordinated with battery energy storage systems (BESS). Their value cannot be inferred from gross wholesale-market revenue: the same dispatch changes retail time-of-use and demand charges, EV user-payment revenue, and service feasibility. This paper develops an execution-aware evaluation framework for a campus EV--BESS portfolio that couples retail bill management with day-ahead (DA), real-time (RT), and ancillary-service (AS) market decisions. The contribution is a physically and financially consistent operating framework, rather than a new MPC or forecasting algorithm. It jointly represents unidirectional workplace EV flexibility, bidirectional BESS dispatch, retail charges, EV service, DA/RT energy settlement, and AS capacity, and evaluates outcomes only from executed dispatch. The RT formulation fixes submitted DA quantities, represents RT quantities as signed increments or decrements, and constrains the actual position as their sum, avoiding an artificial big-M tracking penalty. The framework also clarifies why perfect information over a truncated MPC horizon is not a mathematical upper bound on closed-loop billing-period value. Numerical results in this version are to be regenerated from the corrected formulation before publication.
\end{abstract}

\begin{keyword}
Electric vehicle charging \sep Battery energy storage \sep Model predictive control \sep Value stacking \sep Ancillary services \sep Demand charges \sep Behind-the-meter resources
\end{keyword}

\end{frontmatter}

\section{Introduction}
\label{sec:introduction}

Large EV charging sites are becoming controllable loads at a scale that is relevant for distribution and wholesale-market operation. This trend is reinforced by policy and market changes that seek to lower barriers for distributed energy resources (DERs) to participate in organized wholesale markets through aggregation. In California, distributed resource aggregations can participate in ISO markets for energy and ancillary services subject to metering, telemetry, scheduling-coordinator, and certification requirements \cite{ferc2222}. These developments create a practical question for campuses, commercial sites, and fleet depots: can an existing EV charging portfolio be operated as part of a market-facing DER aggregation without compromising the primary charging service?

The answer is not obvious for workplace charging. Unlike stationary batteries, unidirectional EV chargers cannot export energy from vehicle batteries. Their flexibility is the ability to move charging in time while meeting each driver's departure energy requirement. This flexibility is also intermittent because it exists only when vehicles are plugged in. A BESS can shift energy and hold reserve capability more directly, but behind-the-meter BESS operation also changes the retail meter import that determines time-of-use (TOU) and demand charges. A wholesale dispatch that earns ancillary-service or energy revenue can therefore reduce net value if it creates a new retail demand peak or increases charging during expensive retail periods.

This paper studies that retail--wholesale coupling for a campus EV charging site with a co-located BESS. Its contribution is an execution-aware, market-timeline-consistent operating framework for evaluating whether a workplace portfolio can create net value after wholesale energy and AS participation are layered on top of retail bill management and EV charging revenue. Only the first 15-minute decision from each optimization is implemented; states, commitments, and financial accounts are then updated from the executed quantity. This distinction matters because a planning trajectory can look profitable even if its future actions are later overwritten by MPC updates.

The study is intentionally scoped. The present results use a June--July 2025 UCSD workplace charging case with daytime sessions and essentially no overnight EV carryover. Wholesale prices and retail tariff schedules are treated as known inputs; the EV-information comparison is between perfect EV session information and a persistence EV forecast. Shrinking-horizon and rolling-horizon MPC are both execution structures, rather than competing algorithmic claims. For this workplace setting, shrinking MPC is expected to be a strong benchmark because most flexibility is intraday; rolling MPC is retained for its adaptability to overnight or multi-day settings.

The paper makes four specific contributions:
\begin{enumerate}[leftmargin=*]
    \item It develops an execution-aware campus BTM value-stacking formulation that couples EV charging, BESS dispatch, retail tariffs, EV user payments, DA/RT energy, and AS capacity in one net-value framework.
    \item It treats unidirectional EV charging as a service-constrained virtual battery and coordinates that scheduling flexibility with bidirectional BESS power and SOC headroom under a common meter and market-feasibility envelope.
    \item It introduces a DA/RT settlement-consistent execution layer: the DA energy position remains a fixed financial commitment in RT, while the physical RT position is settled as its deviation rather than being forced by a hard equality or a numerical big-M penalty.
    \item It evaluates the full two-by-two forecast--controller ensemble---perfect versus persistence EV information and shrinking versus rolling horizons---under both retail-only and wholesale-enabled operation, using executed rather than planned trajectories.
\end{enumerate}

The remainder of the paper reviews related EV and DER market-participation literature, defines the campus system and assumptions, presents the EV-information and execution frameworks, gives the optimization formulation, and evaluates the June--July 2025 case study through financial value stacks, daily variability, and representative-day operation.

\section{Related Work}
\label{sec:related_work}

\subsection{DER aggregation and wholesale-market access}
Wholesale-market participation by distributed energy resources (DERs) is increasingly framed as an aggregation problem. FERC Order No.~2222 requires organized wholesale markets to reduce barriers for DER aggregations, and explicitly includes resources such as battery storage, rooftop solar, controllable demand, and EV charging equipment in the DER category \cite{ferc2222}. This policy context motivates campus and commercial sites to ask whether behind-the-meter assets can be aggregated into market-facing flexibility while still satisfying local service obligations.

Recent DER aggregation studies develop market and network mechanisms for representing flexible behind-the-meter resources in wholesale markets. Competitive DERA models derive bid and offer curves for aggregators that combine customer-side generation and consumption \cite{chen2023competitiveDER}. Grid-aware aggregation methods characterize feasible cost curves or operating envelopes so that distribution constraints are not ignored \cite{cui2022gridaware,chen2022dsoDERA}. Generic flexibility-aggregation methods also show how heterogeneous DERs can be represented at the aggregate level for market decisions \cite{muller2017aggregation}. These studies establish that DER aggregations can be formulated as market-facing resources, but they generally do not evaluate a measured campus workplace charging portfolio behind a retail meter with EV user payments, retail demand charges, BESS feasibility, DA/RT energy positions, AS capacity awards, and executed financial accounting in the same case study.

\subsection{EV smart charging and aggregate flexibility}
The smart-charging literature shows that EV charging can be shifted to reduce grid impacts, lower operating cost, or increase system flexibility \cite{clement2010impact,sortomme2011optimal,gan2013optimal,richardson2013electric,mwasilu2014electric}. In this literature, EV flexibility is created by the difference between plug-in time and required charging time; it is limited by arrival time, departure time, requested energy, and charger power. Aggregate-flexibility studies further show that many individual EV constraints can be approximated or aggregated for market and control decisions \cite{alTaha2022aggregateFlex,nazari2022updated}.

The present study uses this idea in a narrower but practically important setting: workplace EV charging without V2G. The EV fleet is not modeled as physical storage and does not export energy. Instead, upward flexibility is load reduction relative to a charging baseline, while downward flexibility is additional charging up to connected-charger limits. This distinction is central to the paper because the BESS provides bidirectional energy shifting, whereas unidirectional EVs provide schedulable demand flexibility.

\subsection{EV aggregators in energy, demand-response, and ancillary-service markets}
Several papers study EV aggregators as market participants. Ancillary-service studies often rely on EV battery discharge or V2G capability and examine whether EV fleets can provide frequency regulation or reserve products while respecting user preferences \cite{clairand2020ancillary,sortomme2012v2gAncillary}. Charging-station market papers analyze strategic bidding, retail pricing, or competition when EV charging stations buy energy from wholesale markets and sell charging service to EV users \cite{mousavi2022evcsWholesale,bayani2021strategic}. These works are useful background, but they differ from the present campus problem because they are commonly EV-only, V2G-oriented, strategic-bidding focused, or not evaluated through a behind-the-meter retail tariff.

Two recent UCSD/TotalEnergies studies are the closest predecessors and must be distinguished carefully. Chen et al.~\cite{chen2024costoptimal} developed a cost-optimal workplace EV aggregator model for demand-response market participation under forecast uncertainty, including demand charge management, TOU energy costs, DA demand-response revenue, EV charging revenue, and 15-minute online feedback control. Chen et al.~\cite{chen2026realtime} extended the EV aggregator formulation by modeling aggregated EV charging as a time-varying battery and using a DA optimizer plus a two-stage RT optimizer to address forecast uncertainty. The present paper builds on this EV-flexibility line but changes the application scope: it adds a co-located BESS, includes DA/RT wholesale energy and AS capacity products, evaluates retail-only versus retail+wholesale operation, and computes value from executed behind-the-meter dispatch rather than from EV-only demand-response scheduling.

\subsection{Behind-the-meter BESS value stacking}
BESS operation has been studied for microgrid scheduling, building demand response, PV self-consumption, and demand-charge management \cite{parisio2014model,wang2018predictive,raoufat2017bess}. Behind the meter, value stacking is difficult because one dispatch action can affect several economic accounts at once. Charging the battery can create future discharge capability or downward reserve capability, but it can also increase retail import and demand-charge exposure. Discharging can reduce a retail peak, but it can also reduce state-of-charge headroom for upward reserve or future energy arbitrage.

This coupling becomes sharper when EV charging is added. EVs and the BESS share the same retail meter, and their flexibility must be allocated among EV service, retail bill management, wholesale energy positions, and AS capacity. Therefore, the relevant performance metric is not BESS revenue, EV charging revenue, or wholesale revenue alone, but net value after retail charges and EV service are accounted for.

\subsection{Research gap and novelty of this paper}
The existing literature supports the individual ingredients of this work: DER aggregation, EV smart charging, EV aggregator market participation, BESS value stacking, and MPC execution. The gap is their execution-aware combination for a campus behind-the-meter portfolio in which unidirectional workplace EV charging and bidirectional BESS dispatch are co-optimized under both retail and wholesale value streams.

The novelty of this paper is therefore applied and accounting-focused rather than algorithmic. It evaluates whether a real campus EV--BESS site can create net value when DA/RT energy and AS capacity participation are layered on top of EV charging service and retail bill management. The results are organized to provide evidence for that claim: gross wholesale revenue is compared against changes in TOU charges, demand charges, EV user payments, BESS feasibility, and completed EV service, all using executed MPC dispatch rather than overwritten planning trajectories.

\section{Methods: Campus System, Market Setting, and Assumptions}
\label{sec:system_description}

\subsection{Behind-the-meter campus portfolio}
The modeled site is a campus workplace EV charging system with a co-located BESS behind a retail meter. The operator collects EV user payments, pays retail energy and demand charges, and schedules the BESS and EV charging power. The retail meter observes the net grid import after EV charging and BESS dispatch. In the broader project architecture, building-load and PV data are available, but they are not active decision inputs in the current reported simulations. The current case should therefore be interpreted as an EV--BESS stage of a larger campus DER platform.

The physical resources have distinct roles. EVs are unidirectional controllable demand: they can shift charging in time, but they cannot discharge to the grid. The BESS is a bidirectional storage asset with power and state-of-charge constraints. Both resources affect the same retail meter and both can contribute to the market-facing flexibility envelope. The optimization must therefore decide how much of the available flexibility should be used for EV service, retail bill reduction, wholesale energy positions, and ancillary-service capacity.

\begin{figure*}[!t]
\centering
\MaybeFigure{system_architecture.png}{width=0.88\textwidth}
\caption{Behind-the-meter EV--BESS value-stacking architecture. The active June--July 2025 simulations coordinate workplace EV charging and BESS dispatch through an MPC controller that uses EV information, tariff parameters, wholesale energy prices, ancillary-service prices, BESS state, and DA/RT settlement inputs. PV and building-load streams are available for a separate sensitivity extension but are not included in the reported EV--BESS results.}
\label{fig:system_architecture}
\end{figure*}

\subsection{Market and tariff setting}
The retail layer includes TOU energy charges, a peak-demand charge, and a non-coincident demand charge. The wholesale layer includes day-ahead (DA) and real-time (RT) energy positions and ancillary-service (AS) capacity products. The model uses the market products as an operational abstraction rather than a complete ISO settlement engine. DA quantities are scheduled before operation, RT decisions update dispatch at 15-minute resolution, and AS awards are constrained by EV and BESS deliverability.

This abstraction is sufficient for the main research question: whether adding wholesale participation improves net value for the campus behind-the-meter portfolio. It is not intended to replace a market-registration or settlement study. Certification requirements, telemetry requirements, minimum bid sizes, and scheduling-coordinator arrangements are outside the current optimization and are treated as deployment requirements rather than simulated decision variables.

\subsection{Data sources and preprocessing}
EV session records are converted into 15-minute interval objects. Each session is assigned an arrival interval, departure interval, requested or delivered energy, and charging power limit. Sessions that cannot be represented consistently at 15-minute resolution are filtered or flagged. The implementation uses a maximum charger power of 6.6 kW per active EV session in the aggregate capability calculation. EV user-payment revenue is computed from delivered charging energy and the user charging price.

The EV baseline $B_t^{\EV}$ defines the reference charging trajectory for flexibility accounting. Because EVs do not export energy, their market-facing upward flexibility is modeled as load reduction relative to this baseline, and their downward flexibility is modeled as additional charging above the baseline while vehicles are connected. Baseline construction is therefore a modeling assumption that directly affects the measured amount of EV flexibility.

Wholesale data include DA and RT locational marginal prices and AS prices for regulation up, regulation down, spinning reserve, and non-spinning reserve. Retail data include TOU energy prices and demand-charge rates. All price and tariff inputs are aligned to the same 15-minute simulation index used for EV and BESS control.

\subsection{Modeling assumptions and scope}
\label{sec:assumptions}
All assumptions used in the current study are collected here so that the scope is explicit.
\begin{enumerate}[leftmargin=*]
    \item EV charging is unidirectional. EV batteries are not discharged to the site or grid, and V2G is not modeled.
    \item The BESS is the only bidirectional energy-storage device. BESS degradation is not modeled explicitly; power, SOC, and daily throughput limits control its operation.
    \item EV service is mandatory. Each modeled session must receive its required energy by departure subject to charger and availability constraints.
    \item Wholesale prices, AS prices, and retail tariff schedules are known exogenous inputs in the present simulations. The paper does not evaluate electricity-price forecasting.
    \item The EV-information comparison is limited to perfect EV information and persistence EV forecasting. Once a vehicle arrives, its realized session information is revealed to the controller.
    \item Active DA energy quantities are fixed financial commitments in RT. The RT physical position is optimized subject to updated feasibility and its deviation from the DA commitment is settled at the RT price; it is neither forced by a hard equality nor controlled by a big-M tracking penalty.
    \item AS awards are constrained by power headroom and BESS duration requirements. Expected deployment factors are used for energy accounting; actual regulation mileage and full settlement rules are not simulated.
    \item Building load and PV are available 15-minute exogenous data streams but are not active inputs in the current EV--BESS results. Their impact will be evaluated as a sensitivity extension using the same controller and settlement logic.
    \item The case study covers June--July 2025 UCSD workplace charging at 15-minute resolution. The tested period contains essentially no overnight EV sessions, so conclusions about rolling-horizon value are limited to daytime workplace charging.
\end{enumerate}

\subsection{Workflow overview}
Figure~\ref{fig:workflow} summarizes the data-to-decision workflow. EV sessions, baseline construction, tariff inputs, market prices, BESS state, and DA/RT references are assembled into the optimization problem. At each RT step, the controller solves over the active horizon, implements only the current 15-minute decision, updates EV and BESS states, and records executed financial quantities. The workflow is designed to ensure that monthly value is computed from actions that were actually executed rather than from overwritten future MPC plans.

\begin{figure*}[!t]
\centering
\MaybeFigure{market_mpc_workflow.png}{width=0.90\textwidth}
\caption{Forecasting and MPC workflow. EV information, tariff inputs, wholesale prices, BESS state, and DA/RT references are converted into an optimization problem. The RT controller updates every 15 minutes, executes the first decision, and records executed quantities for financial accounting.}
\label{fig:workflow}
\end{figure*}

\section{EV Information Scenarios}
The controller requires information about current and future EV sessions because plug-in availability and energy requirements determine the feasible charging envelope. The present study does not compare electricity-price forecasts. Retail tariff schedules, DA/RT energy prices, and AS prices are read from processed data and treated as known inputs. The forecast comparison is therefore an EV-information comparison: how much value is lost when future EV session information is estimated rather than known.

\subsection{Forecast objects}
At each control time $t$, the EV forecast object describes a set of currently connected and future EV sessions over the prediction horizon $\calH_t$. For an EV $i$, the relevant information includes its arrival interval $a_i$, departure interval $d_i$, requested energy $E_i^{\req}$, maximum charging power $\bar p_i^{\EV}$, and availability indicator $b_{t,i}^{\EV}$. Once an EV has arrived, the controller observes its realized session information and uses that information in subsequent optimizations. For EVs that have not yet arrived, the controller either uses perfect information for benchmarking or persistence-based estimates for an implementable forecast.

This distinction is important. In this study, ``perfect'' does not mean perfect electricity-price prediction; the price series are not forecasted. Instead, perfect forecasting means that future EV session information is made available to the optimizer as an information benchmark. Persistence forecasting means that future EV charging demand is estimated from recent historical charging behavior and then translated into future session or aggregate charging requirements.

\subsection{Perfect EV-information benchmark}
The perfect forecast case uses the realized EV session set for the study period when constructing the EV constraints. It answers the question: how well could the same optimization and execution framework perform if future EV arrivals, departures, and energy requests were known? This case is not claimed to be implementable in practice. It is an information benchmark that separates the effect of forecast uncertainty from the effect of the controller structure. A large gap between perfect and persistence cases indicates that improved EV forecasting may have economic value. A small gap indicates that the controller is less sensitive to the forecast errors present in the tested month.

Because perfect EV information is available over the active horizon, shrinking MPC can plan a full service-boundary or remaining-boundary schedule with fewer forecast revisions. Rolling MPC can also use perfect information. This is an information benchmark, not a guaranteed upper bound on closed-loop value; the distinction is proved in Section~\ref{sec:perfect_non_dominance}.

\subsection{Persistence EV forecast}
Persistence forecasting is the main implementable baseline. The idea is that recent or representative historical charging behavior is used to estimate future workplace charging demand. In the current EV forecasting pipeline, two quantities are central: the number of EVs expected to be present and the amount of session energy expected to be requested. These estimates are converted into EV availability and service requirements over the prediction horizon. As time advances, newly arrived EVs reveal their realized information, and the forecast object is updated.

The persistence forecast is intentionally simple. Its role is not to be the final forecasting method for all future deployment cases, but to provide a transparent benchmark that can be reproduced and compared against perfect information. For a campus workplace charging site, persistence can perform reasonably because weekday charging patterns are often recurring. However, it can still misrepresent unusual days, holiday effects, or changes in arrival intensity. Those forecast errors affect the EV flexibility envelope and, through the BTM balance, can also affect BESS dispatch, retail charges, and wholesale-market participation.

\subsection{Forecast construction as an optimization input}
For both perfect and persistence cases, the output of the forecast module is not only a scalar load profile. It is a structured set of EV constraints. The forecast determines which EVs are assumed to be available, how much energy must be delivered before departure, and what aggregate charging capability exists at each 15-minute interval. The controller then combines that EV forecast with BESS state, retail tariff parameters, known market prices, and previous market commitments. In this way, the forecast module determines feasible flexibility, while the optimization module decides how to allocate that flexibility among EV service, retail-cost reduction, and wholesale-market products.
\section{MPC Control Frameworks, Bidding Timeline, and Execution Logic}
The same EV/BESS value-stacking model is solved under two execution frameworks: shrinking-horizon MPC and rolling-horizon MPC. Both frameworks use the same device constraints, tariff accounting, wholesale products, and financial evaluation. They differ in how the horizon is updated and how future decisions are recomputed.

\subsection{Why shrinking and rolling MPC are both studied}
Shrinking-horizon MPC is a natural baseline for the June--July workplace charging case. At the beginning of a daily operating problem, the controller optimizes over the remaining intervals in the day. As the day progresses, the remaining horizon shrinks. This is well matched to daytime workplace charging because most sessions arrive and depart during the same day, and the current data set contains essentially no overnight sessions. For this reason, shrinking MPC is not expected to be weaker than rolling MPC in the present scenario. In several cases it can be stronger because it avoids repeatedly expanding the planning window into periods that contain little additional useful EV flexibility.

Rolling-horizon MPC solves a moving-window problem. After each 15-minute execution step, the controller advances the horizon and rebuilds the forecast object. Rolling MPC is therefore more general. It is expected to become more valuable when overnight charging sessions, residential charging, fleet depots, or multi-day uncertainty are included. Those settings create cross-day coupling and make repeated forecast updates more valuable. The present paper therefore does not present rolling MPC as a guaranteed improvement for the June--July workplace case. Instead, it compares shrinking and rolling as two execution structures under a common model, and it interprets the results in the context of the available EV flexibility.

\subsection{Market bidding and MPC timeline}
The physical controller must also respect the market-information timeline. Fig.~\ref{fig:bidding_timeline} summarizes the intended sequence. Historical EV/PV/building data are collected before the operating day. In the current reported model, the EV stream is active while PV and building-load streams are retained for later integration. A day-ahead forecast is created on Day--1 and used to solve the DA optimization. The resulting DA wholesale bid is submitted before the DA market deadline. After the DA market clears, awarded DA quantities become commitments or references for the RT controller. During Day 0, the RT controller updates every 15 minutes using realized EV arrivals, updated BESS SOC, known price streams, and the remaining DA commitments. It then solves an RT problem and executes only the current 15-minute dispatch.

\begin{figure*}[!t]
\centering
\MaybeFigure{market_bidding_mpc_timeline.png}{width=0.94\textwidth}
\caption{Market bidding and MPC execution timeline. The DA stage uses Day--1 information to create a forecast, solve the DA optimization, and submit DA energy/AS bids. After DA market clearing, awarded DA quantities are passed to the RT controller. On Day 0, EV data and forecasts are updated every 15 minutes, the RT MPC solves the current dispatch problem, and only the first 15-minute decision is executed.}
\label{fig:bidding_timeline}
\end{figure*}

The DA/RT distinction matters because the RT problem should not freely erase DA quantities that have already been submitted or awarded. The current implementation uses a two-settlement representation rather than a hard equality or an artificial reference-tracking penalty. The active DA energy position is preserved as a fixed financial commitment. RT optimization selects the physically feasible net market position after updated EV availability and BESS SOC are observed, and the difference from the DA commitment is settled at the RT price. This preserves the economic consequence of a DA commitment while avoiding infeasibility caused by imposing a stale physical schedule on newly revealed states.

\subsection{Day-ahead and real-time roles}
The DA optimization determines the schedule and submitted wholesale position for the next operating day. It uses the selected EV forecast, known tariff parameters, known DA price series, and AS price series. The RT optimization is executed at 15-minute resolution. It observes realized states, updates the EV forecast, retains the active DA energy commitment as a fixed settlement parameter, and computes the immediate EV/BESS dispatch and RT deviation quantity. Financial results are computed from executed decisions rather than from future decisions that may later be overwritten.

\subsection{Example: execution at 9:00 AM and 9:15 AM}
Consider Day 0 at 9:00 AM. At this time, the controller knows the EVs that have already arrived and observes their realized session information. It also knows the current BESS SOC and accumulated retail demand-charge states. The EV forecast module fills in future EV sessions over the remaining horizon using either perfect EV information or persistence estimates. The price inputs are read from the processed market and tariff data, not forecasted by the EV module. The MPC then solves the MILP with EV service constraints, BESS SOC constraints, BTM meter balance, retail-cost terms, wholesale revenue terms, AS capability constraints, and active DA settlement parameters.

Only the 9:00--9:15 AM EV and BESS decisions are implemented. The rest of the optimized trajectory is a plan, not an executed schedule. At 9:15 AM, delivered EV energy and BESS SOC are updated, any newly arrived EVs reveal their session information, and the forecast object is rebuilt. Shrinking MPC solves over the remaining part of the day, whereas rolling MPC shifts its moving horizon forward. The cycle repeats until the operating horizon ends.

\begin{table*}[!t]
\centering
\caption{Execution logic for the 9:00 AM to 9:15 AM update. The same physical model is used by both MPC frameworks; only the horizon update and forecast replacement rule differ.}
\label{tab:execution_example}
\scriptsize
\begin{tabularx}{\textwidth}{@{}p{0.16\textwidth}X X@{}}
\toprule
Step & Shrinking-horizon MPC & Rolling-horizon MPC \\
\midrule
At 9:00 AM & Build a forecast and optimize over the remaining intervals of Day 0. & Build a forecast and optimize over the rolling window that begins at 9:00 AM.\\
EV information & Arrived EVs use realized session information; future EVs use perfect or persistence EV information. & Same EV information rule, but the forecast window shifts forward at the next step.\\
Market information & DA commitments and known price streams enter the RT problem; the RT deviation is settled at the RT price. & Same DA/RT settlement logic.\\
Execution & Implement only the 9:00--9:15 AM decisions. & Implement only the 9:00--9:15 AM decisions.\\
At 9:15 AM & Remove the executed interval and optimize the remaining day. & Shift the window to start at 9:15 AM and rebuild the forecast object.\\
Expected role & Strong baseline for daytime workplace charging without overnight carryover. & More adaptable for future overnight, multi-day, or higher-uncertainty EV settings.\\
\bottomrule
\end{tabularx}
\end{table*}

\subsection{Scenario definitions}
The experimental design is a two-by-two controller and EV-information comparison, crossed with retail-only and wholesale-enabled operation. The controller dimension compares shrinking-horizon MPC with rolling-horizon MPC; the EV-information dimension compares perfect EV information with a persistence EV forecast; and the operating-mode dimension compares retail-only scheduling with retail+wholesale value stacking. Retail-only cases disable wholesale energy and AS participation. Retail+wholesale cases allow DA/RT energy and AS products subject to the same EV/BESS physical constraints and the same retail meter accounting. Because all cases use the same EV service requirement and financial accounting, differences in net value can be attributed to controller execution, EV information, and wholesale participation rather than to inconsistent service delivery.
\section{Optimization Formulation}
This section presents the core mathematical model in the main text. Detailed linearization, binary direction variables, AS duration constraints, and revenue-decomposition equations are provided in Appendix~A. The notation is introduced with the model components rather than placed in a separate nomenclature section, so that each symbol is defined close to where it is used.

The model is solved repeatedly under either shrinking-horizon or rolling-horizon MPC. Let $\calH_t$ denote the active prediction horizon at control time $t$. For shrinking MPC, $\calH_t$ contains the remaining intervals in the operating day. For rolling MPC, $\calH_t$ is a moving look-ahead window that shifts forward after each execution step. Both controllers use the same physical constraints, retail-meter accounting, wholesale-market variables, and EV-service requirements; they differ in how the horizon is updated and how future EV information is replaced as time advances. The controller therefore solves the value-stacking model below at each decision time, executes only the current interval, updates the realized state, and then repeats the process with the appropriate shrinking or rolling horizon.

\begin{table*}[!t]
\centering
\caption{Core notation used in the formulation. Additional product-specific variables are defined where they appear in Appendix~A.}
\label{tab:notation_core}
\scriptsize
\begin{tabularx}{\textwidth}{@{}p{0.16\textwidth}p{0.78\textwidth}@{}}
\toprule
Symbol & Meaning \\
\midrule
$t\in\calT$, $i\in\calI$ & 15-minute interval and EV/session index.\\
$\DeltaT$ & Interval length in hours.\\
$p_{t,i}^{\EV}$, $p_t^{\EV}$ & Charging power of EV $i$ and aggregate EV fleet charging power.\\
$p_t^{\ch,\EV,\WM}$, $p_t^{\dch,\EV,\WM}$, $p_t^{\ch,\EV,\NWM}$, $p_t^{\dch,\EV,\NWM}$ & EV charging/load-reduction components assigned to wholesale (WM) and non-wholesale (NWM) channels.\\
$e_{t,i}$, $E_i^{\req}$ & Delivered energy state and required departure energy of EV $i$.\\
$b_{t,i}^{\EV}$, $p_t^{\EV,\avail}$ & EV plug-in availability indicator and aggregate available charging capability.\\
$p_t^{\ch,\BESS}$, $p_t^{\dch,\BESS}$, $\soc_t^{\BESS}$ & BESS charge power, discharge power, and state of charge.\\
$p_t^{\GI}$ & Behind-the-meter grid import measured at the retail meter.\\
$C_t^{\TOU}$, $C^{\PD}$, $C^{\NCD}$ & Retail energy, peak-demand, and non-coincident demand-charge rates.\\
$C_t^{\LMP,\DA}$, $C_t^{\LMP,\RT}$ & Known DA and RT wholesale energy price inputs.\\
$p_t^{\DA}$, $p_t^{\RT}$, $p_t^{\mathrm{actual}}$ & Fixed DA energy position, signed RT increment/decrement, and actual wholesale energy position, with $p_t^{\mathrm{actual}}=p_t^{\DA}+p_t^{\RT}$.\\
$c_t^{k,\DA}$, $c_t^{k,\RT}$, $c_t^{k,\mathrm{actual}}$ & Fixed DA AS award, signed RT increment/decrement, and actual nonnegative AS quantity for product $k$.\\
$M_{t}^{q,\mathrm{hist}}$ & Executed billing-period demand threshold carried into the current MPC call, for $q\in\{\PD,\NCD\}$.\\
$R^{\EV}$, $R^{\WM}$ & EV user-payment revenue and wholesale-market revenue.\\
\bottomrule
\end{tabularx}
\end{table*}

\subsection{Objective function}
At each optimization call, the controller minimizes net cost, equivalently maximizing net value. The compact objective is
\begin{align}
\min \quad
J = C^{\TOU}+C^{\PD}+C^{\NCD}-R^{\EV}-R^{\WM},
\label{eq:objective_main}
\end{align}
where $C^{\TOU}$ is the retail energy charge, $C^{\PD}$ is the peak-demand charge, $C^{\NCD}$ is the non-coincident demand charge, $R^{\EV}$ is EV user-payment revenue, and $R^{\WM}$ is wholesale revenue under the DA/RT settlement representation. This objective makes the main economic point of the paper visible: wholesale revenue appears with a negative sign in a cost-minimization objective, but it is not free. The same dispatch also changes retail import and therefore retail cost.

\subsection{Expanded objective components}
For reproducibility, the compact objective in Eq.~\eqref{eq:objective_main} can be expanded into interval and billing-period components. The interval EV revenue is
\begin{align}
r_t^{\EV}=C^{\EV}p_t^{\EV}\DeltaT.
\label{eq:interval_ev_rev}
\end{align}
The interval TOU energy cost is
\begin{align}
c_t^{\TOU}=C_t^{\TOU}p_t^{\GI}\DeltaT.
\label{eq:interval_tou}
\end{align}
The wholesale energy and AS revenue can be written as
\begin{align}
r_t^{\WM}=&\DeltaT\left(C_t^{\mathrm{LMP},\DA}p_t^{\DA}+C_t^{\mathrm{LMP},\RT}p_t^{\RT}\right)\nonumber\\
&+\DeltaT\sum_{k\in\calK}\left(C_t^{k,\DA}c_t^{k,\DA}+C_t^{k,\RT}c_t^{k,\RT}\right)+r_t^{\mathrm{deploy}}.
\label{eq:interval_wm_rev}
\end{align}
Demand charges are not simple interval sums. The implemented MPC objective charges only the increase beyond the already executed billing-period threshold. Define the horizon peaks
\begin{align}
z^{\PD}&\ge p_t^{\GI},\quad t\in\calH_t\cap\calT_{\PD},\label{eq:pd_epigraph}\\
z^{\NCD}&\ge p_t^{\GI},\quad t\in\calH_t,\label{eq:ncd_epigraph}\\
C_t^{\PD}&=C^{\PD,rate}\pos{z^{\PD}-M_t^{\PD,\mathrm{hist}}},\nonumber\\
C_t^{\NCD}&=C^{\NCD,rate}\pos{z^{\NCD}-M_t^{\NCD,\mathrm{hist}}}.\label{eq:dem_epigraph_cost}
\end{align}
The positive-part operator is linearized by nonnegative auxiliary variables. After executing the first interval, each threshold is updated from the executed meter power only. This prevents charging the sunk historical maximum again at every MPC solve.

\subsection{EV charging model}
Each EV session has an arrival time $a_i$, departure time $d_i$, requested energy $e_i^{\req}$, and availability indicator $b_{t,i}^{\EV}$. The charging power is constrained by availability and charger rating:
\begin{align}
0 \le p_{t,i}^{\EV} \le \hat P^{\EV} b_{t,i}^{\EV}.
\label{eq:ev_power_bound_main}
\end{align}
The delivered energy state evolves as
\begin{align}
e_{t,i}=e_{t-1,i}+p_{t,i}^{\EV}\DeltaT/\eta_i^{\EV},
\label{eq:ev_energy_main}
\end{align}
and the service requirement is
\begin{align}
e_{d_i,i}\ge e_i^{\req}.
\label{eq:ev_service_main}
\end{align}
The aggregate EV charging power is
\begin{align}
p_t^{\EV}=\sum_{i\in\calV} p_{t,i}^{\EV}.
\label{eq:ev_aggregate_main}
\end{align}

Equation~\eqref{eq:ev_aggregate_main} is more than notation. It turns many individual charging decisions into the fleet-level demand seen by the meter and the wholesale module. Because V2G is not modeled, $p_{t,i}^{\EV}$ and $p_t^{\EV}$ are nonnegative. The EV fleet still has economic flexibility because the optimizer can decide when to deliver the required energy. If the fleet would normally charge according to a baseline $B_t^{\EV}$, then reducing $p_t^{\EV}$ below that baseline releases upward flexibility; increasing $p_t^{\EV}$ above that baseline absorbs additional power and can provide downward flexibility. The flexibility is real only if the final energy requirement \eqref{eq:ev_service_main} remains satisfied.

The dynamic aggregate charging capability is
\begin{align}
P_{t,\max}^{\EV}=\hat P^{\EV}\sum_{i\in\calV}b_{t,i}^{\EV}.
\label{eq:ev_capability_main}
\end{align}
This value changes throughout the day because vehicles arrive and depart. It determines how much EV demand can be increased at time $t$. It also appears in the AS and net-power envelopes in Appendix~\ref{app:complete_milp}. A high $P_{t,\max}^{\EV}$ does not automatically mean high flexibility: if the connected vehicles are already nearly fully charged, downward flexibility may be limited by remaining energy need; if they are close to departure, upward flexibility may be limited by service completion.

\subsection{BESS model}
The BESS has charging power $p_t^{\ch,\BESS}$, discharging power $p_t^{\dch,\BESS}$, and SOC $\soc_t^{\BESS}$. The SOC dynamics are
\begin{align}
\soc_t^{\BESS}=\soc_{t-1}^{\BESS}+\frac{\DeltaT}{C^{\BESS}}
\left(p_t^{\ch,\BESS}\sqrt{\gamma}-\frac{p_t^{\dch,\BESS}}{\sqrt{\gamma}}\right),
\label{eq:bess_soc_main}
\end{align}
with bounds
\begin{align}
\soc_{\min}^{\BESS}\le \soc_t^{\BESS}\le \soc_{\max}^{\BESS}.
\label{eq:bess_soc_bound_main}
\end{align}
The daily terminal equality is
\begin{align}
\soc_{T_d}^{\BESS}=\soc_{0,d}^{\BESS},
\label{eq:bess_terminal_main}
\end{align}
where $T_d$ denotes the end of the daily accounting horizon. This prevents the controller from artificially improving a daily objective by emptying or filling the BESS at the end of the simulated day.

The BESS is the physical resource that can actually move energy across time. It can charge when wholesale or retail conditions favor absorption, discharge when peak reduction or energy revenue is valuable, and hold SOC to provide upward or downward AS capability. It therefore complements the EV fleet. EV flexibility is constrained by user service and plug-in availability; BESS flexibility is constrained by SOC, power, and throughput.

\subsection{Behind-the-meter power balance}
The retail meter import is modeled as
\begin{align}
p_t^{\GI}=p_t^{\load}+p_t^{\EV}+p_t^{\ch,\BESS}-p_t^{\dch,\BESS}-p_t^{\PV}.
\label{eq:btm_balance_main}
\end{align}
In the reported EV--BESS results, $p_t^{\load}$ and $p_t^{\PV}$ are set to zero. Equation~\eqref{eq:btm_balance_main} is retained because it is the general BTM balance and directly defines the planned PV/building sensitivity extension. In that extension, building load and PV are exogenous, time-varying parameters rather than decision variables; the controller continues to optimize only EV charging, BESS operation, wholesale positions, and AS awards. Perfect and persistence forecasts can be applied to these two exogenous streams without changing the MILP structure.

This balance is the bridge between wholesale actions and retail costs. If the BESS charges to support regulation down or energy arbitrage, $p_t^{\GI}$ can increase. If that happens during an expensive TOU period or during the monthly peak, retail costs can rise. Likewise, if EV charging is shifted to avoid a demand charge, the site may lose wholesale opportunities. This is why the model optimizes net value rather than maximizing $R^{\WM}$ alone.

\subsection{Retail tariff terms}
The TOU energy charge is
\begin{align}
C^{\TOU}=\sum_{t\in\calT} C_t^{\TOU}p_t^{\GI}\DeltaT.
\label{eq:tou_main}
\end{align}
The meter variable is signed. Hence $p_t^{\GI}>0$ is import and $p_t^{\GI}<0$ is export. The adopted tariff convention intentionally settles both signs at $C_t^{\TOU}$; no separate import/export variables or prices are introduced.

Let $M_k^{q,\mathrm{hist}}$ be the maximum relevant executed import before MPC call $k$, and let $z_k^q$ be the maximum predicted import inside the active horizon, for $q\in\{\PD,\NCD\}$. The implemented incremental demand-charge objective is
\begin{align}
C_k^{q,\mathrm{inc}}=C^{q,\mathrm{rate}}\pos{z_k^q-M_k^{q,\mathrm{hist}}},
\label{eq:demand_charge_main}
\end{align}
where $z_k^{\PD}=\max_{t\in\calH_k\cap\calT_{\PD}}p_t^{\GI}$ and $z_k^{\NCD}=\max_{t\in\calH_k}p_t^{\GI}$. After the first decision is executed,
\begin{align}
M_{k+1}^{q,\mathrm{hist}}=max\{M_k^{q,\mathrm{hist}},\,m_q(k)p_k^{\GI,\executed}\},
\label{eq:threshold_update}
\end{align}
where $m_q(k)$ is the tariff eligibility indicator. This state update makes the billing-period maximum path dependent.

\subsection{Wholesale revenue and AS products}
Define the signed RT energy adjustment and actual energy position by
\begin{align}
p_t^{\mathrm{actual}}=p_t^{\DA}+p_t^{\RT},\qquad p_t^{\RT}\in\mathbb{R}.
\label{eq:pactual_main}
\end{align}
For each AS product $x\in\{\RU,\RD,\SP,\NSP\}$,
\begin{align}
c_t^{x,\mathrm{actual}}=c_t^{x,\DA}+c_t^{x,\RT},\quad
c_t^{x,\DA}\ge0,\quad c_t^{x,\RT}\in\mathbb{R},\quad c_t^{x,\mathrm{actual}}\ge0.
\label{eq:cactual_main}
\end{align}
Thus an RT variable is an increment or decrement, not a second nonnegative total award. Wholesale revenue is
\begin{align}
R^{\WM}=&\sum_{t\in\calT}\DeltaT\left(C_t^{\mathrm{LMP},\DA}p_t^{\DA}+C_t^{\mathrm{LMP},\RT}p_t^{\RT}\right)\nonumber\\
&+\sum_{t\in\calT}\DeltaT\sum_{x\in\calK}
\left(C_t^{x,\DA}c_t^{x,\DA}+C_t^{x,\RT}c_t^{x,\RT}\right)+R^{\mathrm{deploy}},
\label{eq:wm_revenue_main}
\end{align}
where expected activation is priced at the RT energy price:
\begin{align}
R^{\mathrm{deploy}}=\sum_t\DeltaT C_t^{\mathrm{LMP},\RT}
\left(\alpha_{\RU}c_t^{\RU,\mathrm{actual}}-\alpha_{\RD}c_t^{\RD,\mathrm{actual}}+\alpha_{\SP}c_t^{\SP,\mathrm{actual}}+\alpha_{\NSP}c_t^{\NSP,\mathrm{actual}}\right).
\label{eq:deployment_revenue_main}
\end{align}

The capacity variables $c_t^{k,s}$ are not unconstrained bids. They are bounded by EV and BESS power headroom, BESS SOC duration requirements, EV baseline deviations, and directionality constraints. For example, the combined upward capability includes BESS discharge headroom plus the ability to reduce EV charging relative to baseline. The combined downward capability includes BESS charge headroom plus the ability to increase EV charging up to $P_{t,\max}^{\EV}$. The complete envelopes are listed in Appendix~\ref{app:complete_milp}.

\subsection{Why EV flexibility and BESS dispatch must be co-optimized}
A tempting simplification is to first optimize EV charging for user service and then optimize the BESS for market value. That decomposition is not valid for this paper's value-stacking problem because EV charging and BESS dispatch share the same meter, the same retail demand-charge state, and the same wholesale feasibility envelope. If EV charging is delayed to reduce meter import during a peak-demand interval, the BESS may have more opportunity to charge or may need to discharge less. If the BESS charges to support a downward product, the EV fleet may need to reduce charging to avoid exceeding a demand-charge threshold. These interactions appear directly in the meter balance \eqref{eq:btm_balance_main} and the capability constraints in Appendix~\ref{app:complete_milp}.

The formulation also avoids double counting EV flexibility. EVs can only be used if the resulting charging schedule still delivers the required energy by departure. In other words, EV flexibility is a scheduling margin, not a new energy source. The model can reduce charging now only if it can charge later before departure. The aggregate baseline and capability variables make this margin visible to the wholesale module, while the individual EV energy constraints protect user service.

\subsection{DA/RT settlement-consistent execution}
The DA and RT problems share the same physical EV/BESS model. Let $m_t^{\DA}\in\{0,1\}$ identify an interval whose submitted DA commitment is active and let $\bar p_t^{\DA}$ be that submitted schedule. In RT,
\begin{align}
p_t^{\DA}=m_t^{\DA}\bar p_t^{\DA}
\label{eq:da_commitment}
\end{align}
is a fixed parameter, while $p_t^{\RT}$ is a free signed adjustment. The physical and financial identities are
\begin{align}
p_t^{\mathrm{actual}}=p_t^{\DA}+p_t^{\RT},\qquad
r_t^{\mathrm{energy}}=\DeltaT\left(C_t^{\mathrm{LMP},\DA}p_t^{\DA}+C_t^{\mathrm{LMP},\RT}p_t^{\RT}\right).
\label{eq:rt_deviation}
\end{align}
The former term $M|p_t^{\mathrm{actual}}-p_t^{\DA}|$ is absent. Following DA is represented economically by the two-settlement cash flow and physically through feasibility of $p_t^{\mathrm{actual}}$; it is not enforced by an artificial tracking penalty or big-M constraint. The same construction is used for AS through Eq.~\eqref{eq:cactual_main}.
\subsection{EV user payments and service economics}
EV user revenue is modeled as
\begin{align}
R^{\EV}=\sum_{t\in\calT} C^{\EV}p_t^{\EV}\DeltaT.
\label{eq:ev_revenue_main}
\end{align}
This term represents user payment for delivered charging energy. It does not mean that the optimizer should overcharge vehicles; session requirements and upper energy bounds prevent that. In the June--July 2025 accounting, EV user revenue is the same across completed cases when the same EV service is delivered. The differences in net value therefore come from retail costs and wholesale revenue rather than from changing total EV energy served.

\subsection{Retail-only mode as an embedded benchmark}
Retail-only operation is not a different physical model. It is the restriction
\begin{align}
p_t^{\DA}=p_t^{\RT}=0,\qquad
c_t^{x,\DA}=c_t^{x,\RT}=0,quad x\in\calK,
\label{eq:retail_only}
\end{align}
together with zero WM allocation variables. EV service, BESS physics, the meter balance, TOU charges, and demand-charge states remain unchanged. Consequently, wholesale-enabled operation weakly enlarges the feasible set of a \emph{single optimization call} whenever zero market participation is feasible. This does not by itself imply that a sequence of receding-horizon calls must dominate another closed-loop policy over an entire billing period.

\subsection{Demand-charge formulations: equivalence and horizon mismatch}
Let $z$ be the maximum meter import predicted inside the active horizon and let $M\ge0$ be the fixed historical threshold. The former total-peak expression and the implemented incremental expression are
\begin{align}
\Phi_{\mathrm{old}}(z;M)&=C^{q,\mathrm{rate}}\max\{z,M\},\label{eq:dc_old}\\
\Phi_{\mathrm{inc}}(z;M)&=C^{q,\mathrm{rate}}\pos{z-M}.\label{eq:dc_inc}
\end{align}
Because
\begin{align}
\Phi_{\mathrm{old}}(z;M)=C^{q,\mathrm{rate}}M+\Phi_{\mathrm{inc}}(z;M),
\label{eq:dc_equivalence}
\end{align}
the two expressions have exactly the same minimizers when $M$ is constant within the solve and both use the same physical meter variable. The difference is accounting: Eq.~\eqref{eq:dc_old} includes a sunk historical bill component, whereas Eq.~\eqref{eq:dc_inc} reports only the avoidable increment. A genuinely different model results if the old expression is applied to EV-only charging rather than $p_t^{\GI}$, because that omits BESS-induced meter peaks.

Neither expression supplies information about peaks beyond the prediction horizon. A future-aware terminal approximation may define a forecast of the remaining billing-period peak $\widehat M_{k}^{q,\mathrm{tail}}$ and an effective threshold
\begin{align}
M_k^{q,\mathrm{eff}}=\max\{M_k^{q,\mathrm{hist}},\widehat M_k^{q,\mathrm{tail}}\},
\label{eq:effective_threshold}
\end{align}
then use
\begin{align}
\widehat C_k^{q,\mathrm{inc}}=C^{q,\mathrm{rate}}\pos{z_k^q-M_k^{q,\mathrm{eff}}}.
\label{eq:future_aware_dc}
\end{align}
This is a proposed extension, not the current code. The current two base MPC notebooks implement Eqs.~\eqref{eq:demand_charge_main}--\eqref{eq:threshold_update}, while the corrected actual-power capability Eqs.~\eqref{eq:pactual_main}, \eqref{eq:app_upenv}, and \eqref{eq:app_downenv} are already implemented. If $\widehat M_k^{q,\mathrm{tail}}$ is too small, the controller overvalues present peak shaving and may spend BESS energy too early. If it is too large, the controller undervalues present peak shaving and may permit an avoidable peak. Suitable estimators include a conservative quantile of forecast net load, a scenario-based expected terminal value, or a learned approximation to the billing-period cost-to-go.

\section{Why Perfect Information Need Not Dominate in Receding-Horizon MPC}
\label{sec:perfect_non_dominance}

\subsection{The dominance theorem for a true full-horizon oracle}
Let $\omega$ denote the complete realized data trajectory over evaluation set $\calT_E$, including EV sessions and all exogenous inputs, and let $\mathcal X(\omega)$ be the corresponding full-horizon feasible set. A genuine perfect-information oracle solves
\begin{align}
J^{\star}(\omega)=\min_{x\in\mathcal X(\omega)}J_{\calT_E}(x;\omega).
\label{eq:oracle_problem}
\end{align}
For every other feasible realized schedule $x^{\pi}\in\mathcal X(\omega)$ generated by policy $\pi$,
\begin{align}
J^{\star}(\omega)\le J_{\calT_E}(x^{\pi};\omega).
\label{eq:oracle_dominance}
\end{align}
Under the revenue convention $V=-J$, this is $V^{\star}(\omega)\ge V^{\pi}(\omega)$. Therefore ``Perfect must be best'' is mathematically correct only when Perfect means one optimization over the complete evaluation horizon with the same objective, constraints, initial state, terminal conditions, and realized data.

\subsection{The implemented Perfect case is an information policy, not that oracle}
At execution time $k$, an MPC policy solves a truncated problem
\begin{align}
x_{k:k+H_k-1}^{\pi}\in\arg\min_x
\left\{\sum_{\tau=k}^{k+H_k-1}\ell_\tau(x_\tau;\mathcal I_k^{\pi})+
\widehat V_{k+H_k}(s_{k+H_k})\right\},
\label{eq:mpc_bellman}
\end{align}
executes only $x_k^{\pi}$, and repeats. Perfect and Persistence change the information set $\mathcal I_k^{\pi}$, but the current implementation does not use the exact continuation value $V_{k+H_k}^{\star}$. Hence the Bellman principle required for global dominance is absent.

The closed-loop ordering can reverse for the following mathematically distinct reasons.

\begin{enumerate}[leftmargin=*]
\item \textbf{Truncation and missing cost-to-go.} A locally optimal first action can leave an inferior future SOC, EV flexibility distribution, or demand threshold. Persistence may accidentally preserve a state that is more valuable outside the current horizon.

\item \textbf{Path-dependent demand charges.} The state contains $M_k^{q,\mathrm{hist}}$. Two policies with the same initial and final SOC can generate different intermediate thresholds, so equality of endpoint SOC does not imply equality of feasible continuation sets or total cost.

\item \textbf{Irreversible market commitments.} DA energy and AS decisions become fixed parameters in RT. Better information inside one truncated DA solve can select a commitment that is locally attractive but less favorable after subsequent receding-horizon updates.

\item \textbf{Information and controller are not separable.} More accurate EV information changes the first-stage action. It does not merely evaluate the same action more accurately. Without an exact terminal value, a better forecast can therefore move the system toward a worse out-of-horizon state.

\item \textbf{Value-stack interaction.} Write $V^{\pi}_{\mathrm{full}}=V^{\pi}_{\mathrm{retail}}+\Delta V^{\pi}_{\mathrm{WM}}$. Even if $V^{P}_{\mathrm{retail}}\ge V^{Q}_{\mathrm{retail}}$, Persistence $Q$ can rank higher in full mode when
\begin{align}
\Delta V^{Q}_{\mathrm{WM}}-\Delta V^{P}_{\mathrm{WM}}
>V^{P}_{\mathrm{retail}}-V^{Q}_{\mathrm{retail}}.
\label{eq:ranking_reversal}
\end{align}
Adding WM enlarges each single-call feasible set, but it need not improve two different closed-loop policies by the same amount.

\item \textbf{Mixed-integer near-optimality.} With solver tolerance $\varepsilon>0$, several solutions can be nearly tied in the current objective but have different first actions. Repeated execution can amplify these differences. This is a secondary numerical mechanism, not the primary structural explanation.
\end{enumerate}

Accordingly, the correct validation rule is: full-horizon oracle Perfect must dominate; receding-horizon Perfect versus Persistence is an empirical policy comparison, not a theorem. Retail-only Perfect equality between shrinking and rolling is expected only when their active feasible sets, boundary conditions, and information sets coincide at every executed step.

\subsection{Shrinking and rolling horizons in symbolic form}
Let $T_b(k)$ be the active service-boundary endpoint and let $H^{\mathrm{roll}}$ be the nominal rolling look-ahead. Then
\begin{align}
\calH_k^{\mathrm{shrink}}&=\{k,\ldots,T_b(k)\},\label{eq:shrink_horizon}\\
\calH_k^{\mathrm{roll}}&=\{k,\ldots,\min(k+H^{\mathrm{roll}}-1,T_E)\}.\label{eq:roll_horizon}
\end{align}
Both execute only the first action. Shrinking closes at a service boundary; rolling can preserve cross-boundary look-ahead and must separately enforce each calendar-period throughput and terminal condition. Neither controller generally contains the complete billing-period demand-charge horizon.

\section{Case Study}
\label{sec:case_study}

\subsection{UCSD workplace charging setting}
The case study represents a UCSD campus workplace charging data set for June--July 2025. The simulation interval is 15 minutes. The reported result set focuses on EV charging and a 250 kW / 332 kWh BESS with SOC bounds of 0.05--0.95 and round-trip efficiency parameter $\gamma=0.9$. The EV charger limit used in the aggregate capability calculation is 6.6 kW per connected EV session. The solver configuration uses CVXPY and Gurobi with a 1\% MIP gap and six solver threads. The current source has no ordinary per-step RT time cutoff; accepted steps must reach the configured optimality target, while a separate watchdog audit is retained to detect abnormal runs. These are implementation settings rather than universal recommendations.

\begin{table*}[!t]
\centering
\caption{Case-study parameters used in the current simulation implementation. Values correspond to the June--July 2025 EV--BESS evaluation.}
\label{tab:case_parameters}
\scriptsize
\begin{tabularx}{\textwidth}{@{}p{0.23\textwidth}p{0.18\textwidth}X@{}}
\toprule
Parameter & Value & Notes \\
\midrule
Simulation period & June--July 2025 & 60-day evaluation period used for the reported results.\\
Time resolution & 15 min & One day has 96 intervals.\\
BESS power rating & 250 kW & Simulation setting.\\
BESS energy capacity & 332 kWh & Simulation setting.\\
BESS SOC bounds & 0.05--0.95 & Daily terminal equality used for daily accounting.\\
BESS efficiency parameter & $\gamma=0.9$ & Implemented with $\sqrt{\gamma}$ charge/discharge split.\\
EV charger rating & 6.6 kW/session & Used for aggregate EV capability.\\
EV charging price & \$0.40/kWh & Used for EV user-payment revenue.\\
Peak-demand rate & \$3.05/kW & Summer tariff setting used in the current simulation.\\
Non-coincident demand rate & \$15.38/kW & Summer tariff setting used in the current simulation.\\
Solver & Gurobi via CVXPY & MILP implementation.\\
MIP gap & $10^{-2}$ & Simulation setting.\\
RT time limit & 1.0 s & Used for repeated RT execution.\\
Solver threads & 6 & Per-optimization Gurobi setting.\\
PV/building load & Not active in reported results & Available 15-minute data; planned sensitivity extension.\\
\bottomrule
\end{tabularx}
\end{table*}

\subsection{EV data processing}
The EV preprocessing converts raw charging-session records into the session-level and interval-level objects required by the optimizer. The session-level data include arrival time, departure time, delivered energy, and driver/session identifiers when available. The interval-level data define which vehicles are connected in each 15-minute period and how much charging power can be assigned. Sessions that cannot be represented consistently at the simulation resolution should be filtered or flagged, because infeasible session data can cause artificial unmet energy or unrealistic flexibility.

The key modeling choice is that the optimizer controls charging power, not session arrival. Once a vehicle is observed, its current availability is known. Future vehicles are forecast. In a perfect forecast, future arrivals and energy requirements are known from the realized data. In a persistence forecast, future arrivals and session energy are inferred from recent patterns. The preprocessing therefore feeds both the realized execution state and the forecast set used by the MPC.

\subsection{EV baseline construction}
The EV baseline $B_t^{\EV}$ is used to define load deviation and flexibility. In the current simulation workflow, the baseline is generated from a baseline simulation and converted from interval energy to power. This baseline enters the EV net charge/discharge decomposition in Appendix~\ref{app:complete_milp}. It is especially important for AS products, because load reduction or load increase is measured relative to an expected charging trajectory.

The baseline is not merely cosmetic. A baseline determines how much upward or downward flexibility the EV fleet is allowed to claim. If baseline charging assumes all connected EVs charge immediately at maximum power, the fleet may appear to provide large upward flexibility by reducing charging. If the baseline is conservative, downward flexibility may be smaller. Because of this sensitivity, the current results are interpreted as baseline-dependent. The implemented baseline is used consistently across scenarios so that comparisons are internally consistent, but the baseline should be tested in future sensitivity analysis.

\subsection{Market-price processing}
The DA/RT LMP and AS price pipelines align market data to the same 15-minute simulation intervals used by EV and BESS control. DA and RT energy prices enter the energy-revenue term. AS prices enter the capacity-revenue term for regulation up, regulation down, spinning reserve, and non-spinning reserve. The implementation also uses product deployment factors to account for expected energy movement associated with awards. The price-processing pipeline creates the arrays used in the optimizer and the daily diagnostics used to interpret representative-day behavior.

\begin{table*}[!t]
\centering
\caption{Data-processing modules used to create the EV, price, market, dispatch, and accounting inputs.}
\label{tab:data_pipeline}
\scriptsize
\begin{tabularx}{\textwidth}{@{}p{0.20\textwidth}p{0.22\textwidth}p{0.22\textwidth}X@{}}
\toprule
Pipeline component & Main inputs & Main outputs & Manuscript relevance \\
\midrule
EV post-processing & Raw charging-session data, timestamps, session energy, charger information & Clean session table, 15-min availability, arrival/departure indices, delivered/requested energy & Defines $a_i$, $d_i$, $e_i^{\req}$, $b_{t,i}^{\EV}$, and the fleet capability envelope.\\
EV baseline generation & Processed sessions and baseline simulation settings & $B_t^{\EV}$ and interval baseline energy/power & Determines upward/downward EV flexibility relative to the baseline.\\
Persistence EV forecast & Recent historical session energy and EV counts & Forecast session-energy and EV-count profiles & Main implementable EV forecast method.\\
Perfect EV forecast & Realized future session information & Information-benchmark EV forecast & Upper-bound/reference forecast case.\\
LMP processing & CAISO DA/RT market files & Interval-aligned $C_t^{\mathrm{LMP},\DA}$ and $C_t^{\mathrm{LMP},\RT}$ & Drives DA/RT energy revenue.\\
AS processing & CAISO regulation and reserve price files & Interval-aligned $C_t^{k,s}$ arrays & Drives AS capacity revenue.\\
MPC implementation & Forecasts, prices, BESS state, EV state, tariff parameters & DA/RT decisions, dispatch logs, per-day solver outputs & Source of monthly and daily result tables.\\
Financial post-processing & Dispatch logs and price/tariff data & Monthly summary, daily detail, product-level WM tables & Source of results section and appendix accounting tables.\\
\bottomrule
\end{tabularx}
\end{table*}

\subsection{Result traceability}
Each quantitative result is linked to an executed dispatch or accounting file. Monthly value stacks are computed from financial post-processing; daily figures are computed from executed daily records; representative-day traces are taken from per-day solver outputs; and EV-service statements are checked against processed session and execution logs. This traceability is important because the paper compares controller and EV-information choices. A difference between cases should reflect the tested market mode, forecast setting, or horizon rule rather than a change in input data.

\section{Financial Accounting and Evaluation Metrics}
\label{sec:metrics}

The financial evaluation is performed on executed decisions. The primary metric is monthly net value:
\begin{align}
V^{\mathrm{month}}=R^{\EV}+R^{\WM}-C^{\TOU}-C^{\PD}-C^{\NCD}.
\label{eq:month_value_metric}
\end{align}
This metric is reported for retail-only and wholesale-enabled modes. Retail-only mode sets wholesale participation to zero and evaluates the EV-BESS controller under retail tariff and EV user-payment terms. Wholesale-enabled mode allows DA/RT energy and AS capacity decisions. Comparing these modes reveals whether wholesale participation improves net value after retail impacts.

Daily metrics are also computed, but they must be interpreted carefully because demand charges are monthly or billing-period quantities. A daily table can assign incremental demand-charge impacts to the day on which a new peak occurs, but the economic meaning is different from an energy charge that accrues independently each interval. For this reason, the paper reports both monthly summaries and daily traces. Monthly summaries provide the final objective comparison; daily traces diagnose when high-impact events occur.

Wholesale revenue is decomposed by product and by asset where possible. The product decomposition separates DA energy, RT energy, regulation up, regulation down, spinning reserve, and non-spinning reserve. The asset decomposition separates BESS and EV contributions using implementation variables and ex-post attribution. This attribution should be interpreted as accounting, not as a separate physical market settlement.

\section{Results}
\label{sec:results}

The results are organized to test the paper's application claim: a campus EV--BESS portfolio should be evaluated by net retail--wholesale value, not by gross wholesale revenue alone. Three questions guide the presentation. First, does wholesale participation improve monthly net value after retail charges and EV user payments are included? Second, which components create or erode value on individual days? Third, does the controller structure matter in a workplace data set whose flexibility is mostly intraday?

\subsection{Scenario set and executed accounting}
The scenario matrix crosses two market modes, two EV-information settings, and two execution structures. Retail-only operation disables wholesale energy and AS products. Retail+wholesale operation enables DA/RT energy positions and AS capacity awards subject to the same EV service, BESS, and retail-meter constraints. Perfect EV information provides an upper benchmark for session information; persistence forecasting represents a simple implementable EV-information baseline. Shrinking and rolling horizons use the same optimization model but update the active horizon differently.

\begin{algorithm}[!t]
\caption{Execution-aware EV--BESS simulation}
\label{alg:mpc_execution}
\begin{algorithmic}[1]
\STATE Select market mode, EV-information setting, and controller structure.
\STATE Initialize BESS SOC, EV energy states, retail demand-charge states, and DA/RT references.
\FOR{each 15-minute interval $t$ in the study month}
    \STATE Reveal arrived EV sessions and update EV energy, BESS SOC, and demand-charge states.
    \STATE Build the EV information set over $\calH_t$ using perfect information or persistence forecasting.
    \STATE Read tariff, DA/RT price, AS price, and active DA-reference inputs over $\calH_t$.
    \STATE Solve the EV--BESS retail--wholesale MILP over the active horizon.
    \STATE Execute only the current 15-minute EV charging and BESS decisions.
    \STATE Record executed retail charges, EV revenue, wholesale revenue, and service metrics.
\ENDFOR
\STATE Aggregate monthly, daily, product-level, and representative-day metrics from executed quantities.
\end{algorithmic}
\end{algorithm}

The last point is important. MPC repeatedly generates future schedules that are later revised. Only the executed first step of each solve is used for financial accounting. This prevents the study from reporting value that existed only in a planning trajectory.

\begin{table*}[!t]
\centering
\caption{Net-value reporting template (US$). ``Both'' enables wholesale DA/RT energy and AS products. Values must be regenerated from the Version~2.6 settlement-consistent model.}
\label{tab:monthly_results}
\scriptsize
\begin{tabular}{@{}llrrr@{}}
\toprule
Controller & EV information & Retail only & Both & Both $-$ retail only \\
\midrule
Shrinking & Perfect & -- & -- & --\\
Shrinking & Persistence & -- & -- & --\\
Rolling & Perfect & -- & -- & --\\
Rolling & Persistence & -- & -- & --\\
\bottomrule
\end{tabular}
\end{table*}

\subsection{Monthly value-stack evidence}
The monthly comparison directly evaluates the retail--wholesale value-stacking claim. A wholesale-enabled case can earn energy and AS revenue, but that revenue is valuable only if it exceeds the induced changes in TOU energy cost, peak-demand charge, and non-coincident demand charge. The monthly value-stack figures therefore report all components together rather than presenting wholesale revenue alone.

\begin{figure*}[!t]
\centering
\MaybeFigure{fig_financial_case_value_stack_subplots.png}{width=0.96\textwidth}
\caption{Executed June--July 2025 financial value stacks. Each panel compares retail-only and wholesale-enabled operation for one controller--forecast pair. Positive bars are EV user revenue and wholesale-market revenue; negative bars are TOU, peak-demand (PD), and non-coincident demand (NCD) costs. The black marker gives the resulting net value.}
\label{fig:value_stack}
\end{figure*}

Table~\ref{tab:monthly_results} is intentionally left as a reporting template until the corrected model is rerun. The required interpretation is based on executed accounting: gross market revenue must be compared with induced TOU, peak-demand, and non-coincident-demand costs. Controller and information rankings must be reported empirically and must not be presented as universal dominance claims.

\subsection{Daily mechanisms behind monthly value}
Monthly totals can hide the operating days that create or erode value. Demand charges are especially discontinuous because one 15-minute interval can set or reset a monthly maximum. The daily figures are therefore used to identify whether net-value differences arise from systematic daily behavior or from a small number of high-impact days.

\begin{figure*}[!t]
\centering
\MaybeFigure{fig_financial_case_comparison_subplots.png}{width=0.96\textwidth}
\caption{Daily executed net values over June--July 2025. Orange bars denote retail-only operation, blue bars denote wholesale-enabled operation, and the black trace is the daily difference. The large negative bars at the beginning of each billing month reflect allocation of monthly demand-charge increments to the intervals that establish a new billing-period maximum; they should not be interpreted as independent daily energy costs.}
\label{fig:daily_revenue}
\end{figure*}

\begin{figure*}[!t]
\centering
\MaybeFigure{fig_financial_delta_heatmap.png}{width=0.92\textwidth}
\caption{Incremental June--July 2025 value of wholesale-enabled operation relative to retail-only operation. Each cell reports ``Both $-$ retail only'' for one controller--forecast combination. The decomposition distinguishes total net-value change, wholesale-market revenue, and changes in TOU, peak-demand, and non-coincident demand costs.}
\label{fig:incremental_value_heatmap}
\end{figure*}

\begin{figure*}[!t]
\centering
\MaybeFigure{fig_financial_violin_daily_total.png}{width=0.96\textwidth}
\caption{Distribution of daily executed net financial value over June--July 2025. Violin widths summarize the distribution and points show individual days for retail-only and wholesale-enabled operation. The long lower tails mainly reflect days on which a billing-period demand-charge maximum is established.}
\label{fig:daily_value_distribution}
\end{figure*}

Figures~\ref{fig:daily_revenue}--\ref{fig:daily_value_distribution} provide evidence for the mechanism rather than only the outcome. The heatmap separates gross wholesale revenue from induced retail-cost changes, while the distribution view distinguishes systematic daily performance from billing-period outliers. EV flexibility and BESS dispatch can create wholesale revenue on some days, but retail exposure can offset it on days when grid import approaches or exceeds a demand-charge threshold. This explains why the same gross market revenue can have different net value depending on when it occurs.

\subsection{Representative-day operation}
A representative operating day is used to connect the financial outcome to physical decisions. The six-panel plot links EV charging, BESS power and SOC, grid import, wholesale energy quantities, AS awards, and value components. The purpose is not to show that the controller is complex; it is to show how a wholesale decision propagates through the behind-the-meter balance.

\begin{figure*}[!t]
\centering
\MaybeFigure{fig_representative_day_6panel.png}{width=0.92\textwidth}
\caption{Representative-day operation trace from the RT solver-output folder. The panels link EV charging, BESS operation, grid import, wholesale energy quantities, AS awards, and financial impacts for one high-information operating day.}
\label{fig:representative_day}
\end{figure*}

The operational interpretation follows the same causal chain as the formulation. A high AS price can make it attractive to reserve BESS headroom or modulate EV charging relative to baseline. That action changes the meter import, which can change TOU and demand-charge exposure. The representative-day plot therefore supports the paper's main claim: market-facing EV--BESS scheduling must be evaluated through the retail meter and EV service constraints.

\subsection{Controller and EV-information interpretation}
The controller comparison should not be read as a general ranking of shrinking and rolling MPC. In this June--July workplace data set, most EV sessions begin and end within the same day. Shrinking MPC is therefore well matched to the daily service structure. Rolling MPC is retained because it is a more general execution structure and is expected to be more useful when overnight sessions, fleet depots, residential charging, or multi-day uncertainty are included. In the current case, better gross wholesale performance from a rolling controller does not automatically imply higher net value because the retail-meter consequences can be larger.

The EV-information comparison has a different interpretation. Perfect information is an information benchmark for knowing future sessions, but only the full-horizon oracle in Eq.~\eqref{eq:oracle_problem} is a guaranteed performance upper bound. Persistence is an implementable baseline that uses recent charging patterns. The closed-loop gap therefore measures the interaction of information, horizon truncation, terminal approximation, demand-charge state, and market commitments; it should not be interpreted as forecast accuracy alone.

\section{Discussion}
\label{sec:discussion}

The results support three practical messages for campus DER aggregation. First, gross wholesale revenue is an incomplete metric. A behind-the-meter site must settle with both the wholesale market and the retail tariff, and the retail demand-charge component can dominate individual operating intervals. Second, unidirectional EV charging can contribute to flexibility, but it should not be described as storage. EV flexibility is a scheduling margin constrained by connection time and service requirements. Third, controller design should match the physical flexibility of the data set. For daytime workplace charging without overnight carryover, a daily shrinking horizon is a strong baseline rather than a weak comparator.

These points also define the paper's novelty relative to prior EV market-participation studies. The value of the work is not that MPC is used, nor that a persistence forecast is proposed. The value is that EV charging, BESS operation, retail tariffs, DA/RT wholesale energy, AS capacity, and executed accounting are placed in one campus-scale application study. This framing makes the result useful to SEGAN readers interested in whether DER aggregation concepts translate into net value for real behind-the-meter charging sites.

\section{Limitations and Future Work}
\label{sec:limitations}

The current study is limited to a June--July 2025 sample period and an EV--BESS implementation stage. Additional months and seasons should be simulated before making annual value claims. The next sensitivity analysis will activate the available 15-minute PV and building-load series as uncontrollable net-load inputs in Eq.~\eqref{eq:btm_balance_main}, first under perfect information and then under a simple persistence forecast. This design preserves the EV--BESS base case and isolates the incremental effect of exogenous behind-the-meter variability. The AS settlement model uses expected deployment factors and does not simulate realized regulation mileage, telemetry qualification, or all ISO settlement details. The EV baseline is a key assumption and should be validated against alternative baseline definitions. Finally, the tested workplace period contains essentially no overnight sessions; future studies should evaluate overnight workplace charging, residential charging, fleet depots, and other settings where rolling-horizon MPC may have a clearer operational advantage.

\section{Conclusion}
\label{sec:conclusion}

This paper evaluates campus EV charging and BESS as a behind-the-meter DER aggregation application under retail--wholesale coupling. The formulation combines EV user payments, retail TOU and demand charges, DA/RT settlement-consistent energy positions, AS capacity products, EV service constraints, and BESS feasibility in a single executed-accounting framework. The June--July 2025 UCSD workplace case shows why wholesale participation cannot be evaluated from gross market revenue alone: the same actions that create wholesale revenue can also change retail import and demand-charge exposure. The study also clarifies that rolling MPC is not expected to dominate in a daytime workplace data set with little overnight flexibility; its main advantage is adaptability to future cross-day scenarios. The main conclusion is therefore operational rather than algorithmic: campus EV--BESS value stacking is possible, but the evidence for its value must be based on net retail--wholesale accounting and completed EV service.

\section*{Acknowledgements and AI-assisted preparation disclosure}
The authors acknowledge the UCSD and TotalEnergies collaboration that supported the development of the campus EV--BESS simulation workflow. During manuscript preparation, the authors used AI-assisted writing tools to help revise language, organize text, and improve clarity. The authors reviewed, edited, and verified the manuscript content and take full responsibility for the final text.

\appendix
\section{Complete MILP Formulation}
\label{app:complete_milp}

This appendix records the implementation-level constraints. They are separated from the main text to keep the main formulation readable while preserving reproducibility.

\subsection{EV availability and energy state}
\begin{align}
b_{t,i}^{\EV} &= \begin{cases}1, & a_i\le t\le d_i,\\0, & \text{otherwise},\end{cases}\label{eq:app_ev_avail}\\
e_{a_i,i} &= e_{a_i},\qquad e_{0,i}=0,\label{eq:app_ev_init}\\
e_{t,i} &= e_{t-1,i}+p_{t,i}^{\EV}\DeltaT/\eta_i^{\EV},\label{eq:app_ev_dyn}\\
0&\le e_{t,i}\le e_i^{\req},\qquad
0\le p_{t,i}^{\EV}\le \hat P^{\EV}b_{t,i}^{\EV},\label{eq:app_ev_bounds}\\
e_{d_i,i}&\ge e_i^{\req}.\label{eq:app_ev_req}
\end{align}
For rolling MPC, service closure constraints are applied so that energy required for currently relevant same-day sessions cannot be pushed beyond the valid service window. The precise mask depends on whether the horizon is a same-day shrinking horizon or a rolling cross-day horizon.

\subsection{BESS SOC, split, and terminal equality}
\begin{align}
p_t^{\ch,\BESS}&=p_t^{\ch,\BESS,\WM}+p_t^{\ch,\BESS,\NWM},\label{eq:app_bess_chsplit}\\
p_t^{\dch,\BESS}&=p_t^{\dch,\BESS,\WM}+p_t^{\dch,\BESS,\NWM},\label{eq:app_bess_dchsplit}\\
p_t^{\BESS}&=p_t^{\ch,\BESS}-p_t^{\dch,\BESS},\label{eq:app_bess_power}\\
\soc_t^{\BESS}&=\soc_{t-1}^{\BESS}+\frac{\DeltaT}{C^{\BESS}}
\left(p_t^{\ch,\BESS}\sqrt{\gamma}-\frac{p_t^{\dch,\BESS}}{\sqrt{\gamma}}\right),\label{eq:app_bess_soc}\\
\soc_{\min}^{\BESS}&\le \soc_t^{\BESS}\le \soc_{\max}^{\BESS},\label{eq:app_bess_bounds}\\
\soc_{T_d}^{\BESS}&=\soc_{0,d}^{\BESS}.\label{eq:app_bess_terminal}
\end{align}
The wholesale/non-wholesale split is used for accounting and attribution. The physical SOC equation uses the total charge and discharge powers. For every calendar accounting block $d$ intersected by the horizon, the throughput guard is
\begin{align}
\DeltaT\sum_{t\in\calH_k\cap\calT_d}
\left(p_t^{\ch,\BESS}+p_t^{\dch,\BESS}\right)
\le 2C^{\BESS}\left(\soc_{\max}^{\BESS}-\soc_{\min}^{\BESS}\right).
\label{eq:app_bess_throughput}
\end{align}
In rolling MPC this constraint is formed separately for each calendar block represented in the moving horizon.

\subsection{Direction binaries}
\begin{align}
0&\le p_t^{\ch,\BESS}\le P_{\max}^{\BESS}u_t^{\ch,\BESS},\label{eq:app_bess_chbin}\\
0&\le p_t^{\dch,\BESS}\le P_{\max}^{\BESS}u_t^{\dch,\BESS},\label{eq:app_bess_dchbin}\\
u_t^{\ch,\BESS}+u_t^{\dch,\BESS}&\le 1.\label{eq:app_bess_dirbin}
\end{align}
Analogous direction binaries are used for EV net load-increase/load-reduction variables and for net wholesale charge/discharge accounting. These binaries prevent the model from claiming simultaneous upward and downward capability from the same physical margin.

\subsection{EV net-deviation and big-M formulation}
The EV fleet is decomposed relative to the baseline into wholesale and non-wholesale channels:
\begin{align}
p_t^{\EV}-B_t^{\EV}
&=\left(p_t^{\ch,\EV,\WM}-p_t^{\dch,\EV,\WM}\right)
+\left(p_t^{\ch,\EV,\NWM}-p_t^{\dch,\EV,\NWM}\right),\label{eq:app_ev_deviation}\\
p_t^{\ch,\EV}&=p_t^{\ch,\EV,\WM}+p_t^{\ch,\EV,\NWM},\qquad
p_t^{\dch,\EV}=p_t^{\dch,\EV,\WM}+p_t^{\dch,\EV,\NWM},\label{eq:app_ev_channel_split}\\
0&\le p_t^{\ch,\EV}\le M_t^{\EV}u_t^{\ch,\EV},\label{eq:app_ev_ch}\\
0&\le p_t^{\dch,\EV}\le M_t^{\EV}u_t^{\dch,\EV},\label{eq:app_ev_dch}\\
u_t^{\ch,\EV}+u_t^{\dch,\EV}&\le 1,\label{eq:app_ev_dir}\\
M_t^{\EV}&=1.05\max\left\{P_{t,\max}^{\EV},\,|B_t^{\EV}|,\,\epsilon\right\}.\label{eq:app_ev_bigM}
\end{align}
Here $p_t^{\dch,\EV}$ is a load-reduction variable, not physical EV discharge. The WM terms contribute to the net market-facing movement in Eq.~\eqref{eq:app_net_wm_flow}; the NWM terms allow the same EV flexibility to be used independently for retail-meter and demand-charge management. Both channels remain subject to one aggregate EV charging schedule and one service-completion requirement, so the split is an accounting and feasibility allocation rather than double-counted flexibility.

\subsection{Wholesale capacity envelopes}
The total upward and downward deliverable capabilities are
\begin{align}
\hat p_t^{\up}&=P_{\max}^{\BESS}+B_t^{\EV},\label{eq:app_upcap}\\
\hat p_t^{\down}&=P_{\max}^{\BESS}-B_t^{\EV}+P_{t,\max}^{\EV}.\label{eq:app_downcap}
\end{align}
The executed net wholesale position and AS awards must satisfy
\begin{align}
&\pos{p_t^{\mathrm{actual}}}
+c_t^{\RU,\mathrm{actual}}+c_t^{\SP,\mathrm{actual}}+c_t^{\NSP,\mathrm{actual}}
\le \hat p_t^{\up},\label{eq:app_upenv}\\
&\pos{-p_t^{\mathrm{actual}}}+c_t^{\RD,\mathrm{actual}}
\le \hat p_t^{\down}.\label{eq:app_downenv}
\end{align}
Positive and negative components are represented with auxiliary variables in the implementation. Crucially, the envelopes use the actual DA-plus-RT energy position, not the RT adjustment alone.

\subsection{Net power flow and AS deployment accounting}
The market-facing net charge/discharge flow is allocated only to the WM portions of the BESS and EV flexibility:
\begin{align}
p_t^{\ch,\net}-p_t^{\dch,\net}
={}&\left(p_t^{\ch,\BESS,\WM}-p_t^{\dch,\BESS,\WM}\right)\nonumber\\
&+\left(p_t^{\ch,\EV,\WM}-p_t^{\dch,\EV,\WM}\right).
\label{eq:app_net_wm_flow}
\end{align}
The NWM portions remain in the physical meter balance and are available for retail-cost and demand-charge management, but they do not create an additional wholesale position.

Expected deployment and actual energy jointly determine the market-facing movement:
\begin{align}
q_t^{\WM}={}&p_t^{\mathrm{actual}}
+\alpha_t^{\RU}c_t^{\RU,\mathrm{actual}}
-\alpha_t^{\RD}c_t^{\RD,\mathrm{actual}}\nonumber\\
&+\alpha_t^{\SP}c_t^{\SP,\mathrm{actual}}
+\alpha_t^{\NSP}c_t^{\NSP,\mathrm{actual}}.\label{eq:app_deploy}
\end{align}
The implementation enforces $p_t^{\dch,\net}=\pos{q_t^{\WM}}$ and $p_t^{\ch,\net}=\pos{-q_t^{\WM}}$ with binary linearization, then links these quantities to the WM portions of EV and BESS flexibility.

\subsection{AS duration constraints}
Capacity awards must be deliverable for the required duration $T_{\mathrm{DU}}$:
\begin{align}
\soc_{t+1}^{\BESS}
&\ge \soc_{\min}^{\BESS}+\frac{T_{\mathrm{DU}}}{C^{\BESS}\sqrt{\gamma}}
\left(c_t^{\RU,\mathrm{actual}}+c_t^{\SP,\mathrm{actual}}+c_t^{\NSP,\mathrm{actual}}\right),\label{eq:app_duration_up}\\
\soc_{t+1}^{\BESS}
&\le \soc_{\max}^{\BESS}-\frac{T_{\mathrm{DU}}\sqrt{\gamma}}{C^{\BESS}}c_t^{\RD,\mathrm{actual}}.\label{eq:app_duration_down}
\end{align}
The actual AS quantities must remain nonnegative. A negative $c_t^{x,\RT}$ is permitted only to reduce a DA award; it cannot create a negative physical service.

\subsection{DA commitment and RT deviation settlement}
The active DA energy schedule is carried into RT through a binary commitment mask $m_t^{\DA}$:
\begin{align}
p_t^{\DA}&=m_t^{\DA}\bar p_t^{\DA},\label{eq:app_da_commitment}\\
p_t^{\mathrm{actual}}&=p_t^{\DA}+p_t^{\RT}.\label{eq:app_rt_deviation}
\end{align}
For every AS product,
\begin{align}
c_t^{x,\DA}=m_t^{x,\DA}\bar c_t^{x,\DA},\qquad
c_t^{x,\mathrm{actual}}=c_t^{x,\DA}+c_t^{x,\RT}\ge0.
\label{eq:app_as_commitment}
\end{align}
The DA quantities are fixed parameters in RT, while the RT quantities are signed increments/decrements. This is the settlement-consistent replacement for the former big-M tracking penalty.

\subsection{Revenue decomposition}
Wholesale revenue is decomposed as
\begin{align}
R^{\WM,\energy}&=\sum_t\DeltaT\left(C_t^{\mathrm{LMP},\DA}p_t^{\DA}+C_t^{\mathrm{LMP},\RT}p_t^{\RT}\right)+R^{\mathrm{deploy}},\label{eq:app_energy_rev}\\
R^{\WM,\capacity}&=\sum_t\DeltaT\sum_{x\in\calK}\left(C_t^{x,\DA}c_t^{x,\DA}+C_t^{x,\RT}c_t^{x,\RT}\right),\label{eq:app_capacity_rev}\\
R^{\WM}&=R^{\WM,\energy}+R^{\WM,\capacity}.\label{eq:app_total_rev}
\end{align}
Asset-level EV/BESS attribution is an ex-post accounting calculation based on EV and BESS wholesale variables. It should not be interpreted as a separate physical constraint.




\bibliographystyle{elsarticle-num}
\bibliography{references_v2.4}

\end{document}
